% ───────────────────────────────────────────────────────────────────
% Introduction
% ───────────────────────────────────────────────────────────────────

\label{sec:intro}

\subsection{Background}

Large language models are increasingly deployed as autonomous agents that
interact with environments, tools, and other agents.  Recent work has shown
that LLM-powered agents can exhibit emergent social behaviours in open-ended
simulation~\cite{park2023generative}, coordinate within multi-agent
frameworks~\cite{wu2023autogen}, and interleave reasoning with
action~\cite{yao2023react}.  A natural extension of this research agenda is
to ask whether LLM agents can engage in structured economic interaction:
bilateral bargaining under private constraints, with real stakes, limited
information, and measurable outcomes.  Open-ended social simulation
demonstrates that LLM agents can sustain coherent dialogue and exhibit
personality-consistent behaviour, but it does not require agents to satisfy
hard economic feasibility constraints, discover prices, or produce outcomes
that can be evaluated against classical economic benchmarks.  Structured
economic interaction imposes a qualitatively different test.

Bilateral bargaining is among the best-characterised settings in economic
theory.  Rubinstein's alternating-offer model~\cite{rubinstein1982} predicts
that patient agents extract larger surplus shares through time-preference
discounting.  The anchoring literature~\cite{tversky1974} establishes anchoring
as a general cognitive heuristic; in negotiation contexts, Galinsky and
Mussweiler~\cite{galinsky2001} show that first offers systematically bias
final agreements.  Deadline pressure effects have been documented in human
experiments~\cite{roth1988deadline}, and decentralised bilateral markets are
expected to produce equilibrium prices with persistent price
dispersion~\cite{rubinstein1985}.  These theoretical predictions provide
reference points against which observed agent behaviour can be compared
directionally, though the extent to which any simulator-based evaluation
constitutes a clean test depends on the prompt and protocol design.

Despite growing interest in LLM negotiation, existing studies typically
focus on multi-issue dialogue games~\cite{lewis2017dealornodeal,
abdelnabi2024llmdeliberation}, personality-driven
interaction~\cite{mao2024personalities}, or assistive human--AI
bargaining~\cite{chawla2023assistive}.  These designs evaluate negotiation
capability but do not systematically test LLM agents against the specific
quantitative predictions of classical bargaining theory, nor do they
examine whether LLM agents reproduce qualitative behavioural patterns
documented in human negotiation studies---such as deadline-driven
concession acceleration~\cite{roth1988deadline}, Boulware-style
tactical hardening, or first-offer anchoring~\cite{galinsky2001}.
Furthermore, most studies operate at a single level of analysis without
connecting micro-level bargaining outcomes to macro-level market dynamics
within a single controlled platform.  The practical relevance of this gap
is growing as automated negotiation enters consumer and enterprise
markets~\cite{davidson2024automatedrisky}, motivating controlled
simulator-based evaluation of LLM bargaining behaviour even where
conclusions remain preliminary.

\subsection{Problem Statement}

The central question motivating this thesis is:

\begin{quote}
\textit{What patterns of concession, price discovery, and market efficiency
emerge when LLM agents negotiate under controlled bilateral bargaining
conditions, and to what extent are observed outcomes shaped by the
informational and protocol environment in which agents operate?}
\end{quote}

This question cannot be answered by evaluating LLM outputs in isolation.
Simulator design choices---the prompt architecture that structures agent
reasoning, the protocol design that governs turn-taking and termination, and
any enforcement logic applied to agent outputs---all materially alter what
is observable.  These design choices must be made explicit and independently
configurable to enable defensible interpretation of experimental findings.
A further challenge is attribution: any observed outcome in an LLM
simulation reflects the combined system of model, prompt, and protocol.
Disentangling the contribution of LLM decision-making from the contribution
of simulator infrastructure requires careful logging, explicit ablation
studies, and conservative interpretation.  This thesis builds a platform
designed to support such analysis, though it does not fully resolve the
attribution problem within the scope of a single project.  The simulator
is the instrument; the primary subject of investigation is the behaviour of
LLM agents within it.

\subsection{Research Objectives}

A configurable multi-agent negotiation simulator was designed and
implemented to support the empirical programme described below.  The system
enforces private economic constraints, supports a structured alternating-offers
protocol, and provides event-level logging sufficient for session and
market-level analysis.  The research objectives are:

\begin{enumerate}
    \item Characterise LLM bargaining behaviour empirically---concession
    dynamics, deadline sensitivity, and deal success---and compare observed
    patterns directionally against predictions from classical bargaining
    theory.

    \item Investigate whether and how LLM-agent markets respond to changes
    in fundamentals (supply--demand shifts, exogenous shocks) and
    informational environments (reference-price anchoring), and assess
    whether price discovery emerges from bilateral negotiation under the
    tested conditions.

    \item Evaluate whether institutional mechanism design (matching
    algorithms) can improve market-level welfare independently of agent
    sophistication.

    \item Identify how simulator design choices---prompt architecture,
    feasibility enforcement, and reference-price framing---function as
    protocol-level design considerations, and report methodological cautions relevant to
    LLM-based economic simulation.
\end{enumerate}

\subsection{Contributions}

This thesis makes four contributions.  All results characterise the
combined system (LLM~+~prompt~+~protocol) under specific conditions; they
should not be generalised across model families or scales without further
ablation and replication.

\textbf{1.\ Empirical characterisation of LLM bargaining behaviour.}
LLM agents produce concave (Boulware-style) concession trajectories with
a 5.65:1 buyer-to-seller concession ratio, and deadline pressure
concentrates settlement in the final rounds (82.4\% of deals close in the
final two rounds under a six-round horizon).  Price stickiness is
substantial: demand and supply elasticities of 0.275 and 0.035
respectively indicate that agents update prices only weakly in response to
changing fundamentals.  Several of these patterns are directionally
consistent with behavioural patterns documented in the human negotiation
literature, though no direct human comparison was conducted and the
underlying mechanisms may differ.

\textbf{2.\ Demonstration that prompt-embedded information is a first-order
determinant of market outcomes.}  An anchor-environment experiment
(Experiment~J) shows that the reference-price signal in the agent prompt
produces a mean transaction price spread of approximately \$30 across anchor
modes under identical fundamentals, and reverses the sign of shock effects
between anchor conditions.  This establishes prompt-level information
framing as a first-order experimental variable in LLM-based market
simulation, with direct implications for the interpretation of all other
experiments.

\textbf{3.\ Identification of a micro--macro capability gap.}  LLM agents
are effective deal-closing mechanisms---78.6\% overall deal rate, 81\%
market liquidity---but weak price-discovery mechanisms: prices exhibit
persistent within-tick dispersion without convergence, and elasticities
remain near zero under the tested model--prompt combination.  This
separation between deal-closing capability and price-discovery capability
distinguishes two dimensions of negotiation competence that are often
conflated in prior evaluation frameworks.

\textbf{4.\ A configurable experimental platform} connecting micro-level
bargaining to macro-level market dynamics, with hard private constraints,
structured event logging, and explicit controls over design choices including
gate status, prompt architecture, reference-price mode, and matching
algorithm.  The platform enables the findings above and provides a
foundation for future ablation and extension.

The remainder of this thesis is organised as follows.
Section~\ref{sec:related} reviews related work in LLM-based simulation, LLM
negotiation, and classical bargaining theory.
Section~\ref{sec:formulation} describes the problem formulation and simulator
architecture.
Section~\ref{sec:experiments} details the methodology and experimental design
for all seven experiments.
Section~\ref{sec:results} presents results and analysis.
Section~\ref{sec:discussion} discusses implications, limitations, and the
role of simulator design choices as experimental variables.
Section~\ref{sec:conclusion} concludes and identifies directions for future
work.
